\documentclass{article}

% --- GESTION DES MARGES DE PAGE ---
\usepackage[a4paper, top=2.5cm, bottom=2.5cm, left=3cm, right=3cm]{geometry}

% --- PRÉAMBULE STANDARD ---
\usepackage[utf8]{inputenc}
\usepackage[T1]{fontenc}
\usepackage{lmodern}
\usepackage[french]{babel}
\usepackage{parskip} % NOUVEAU : Supprime l'indentation et ajoute espace inter-paragraphe
\usepackage{xcolor}
\usepackage{tcolorbox}
\usepackage{listings}
\usepackage{amsmath}
\usepackage{graphicx} % Requis pour inclure des images
\usepackage{amssymb} % Pour les symboles mathématiques comme \subseteq
\usepackage{sectsty} % Pour le style des sections
\usepackage{etoolbox} % Pour les conditions
\usepackage[dvipsnames]{xcolor}
\usepackage{pgfplots} % Le package principal pour les graphiques
\usepackage{enumitem}
\usepackage{diagbox}

\setlist[itemize,1]{label=$\cdot$}

\usetikzlibrary{
    matrix, 
    patterns.meta,
    calc,
    positioning,
    decorations.pathreplacing,
    trees,
    backgrounds
}

% Redéfinir le symbole pour le premier niveau de la liste
\renewcommand{\labelitemi}{$\cdot$}
\renewcommand{\labelitemii}{$\circ$}
\renewcommand{\labelitemiii}{$\ast$}

% Styles de hachures (inchangés)
\tikzset{
  red_hatch/.style={
    pattern={Lines[angle=45, line width=0.8pt, distance=4pt]}, 
    pattern color=red
  },
  blue_hatch/.style={
    pattern={Lines[angle=-45, line width=0.8pt, distance=4pt]}, 
    pattern color=blue
  },
  purple_hatch/.style={
    pattern={Lines[angle=45, line width=0.8pt, distance=4pt]}, 
    pattern color=red,
    postaction={
      pattern={Lines[angle=-45, line width=0.8pt, distance=4pt]}, 
      pattern color=blue
    }
  }
}

% --- BIBLIOTHÈQUES TCOLORBOX ---
\tcbuselibrary{listings, skins, breakable}

% --- GESTION DES LIENS HYPERTEXTE ---
\usepackage[colorlinks=true, linkcolor=black, urlcolor=blue]{hyperref}

% --- SOULIGNER LES TITRES ET SOUS-TITRES ---
\sectionfont{\underline}
\subsectionfont{\underline}
\subsubsectionfont{\underline}

% --- PAGE DE GARDE AMÉLIORÉE ---
\makeatletter
\renewcommand{\maketitle}{%
\begin{titlepage}
\centering
\vspace*{\stretch{1.5}}
{\Huge \bfseries Mes Notes de Lecture\par}
\vspace{0.4cm}
\rule{0.8\linewidth}{0.4pt}
\vspace{1cm}
\vspace*{\stretch{2.5}}
\end{titlepage}
}
\makeatother

% ==================================================================
% --- MODIFIÉ : DÉFINITION DES COULEURS STYLE VS CODE (LIGHT) ---
% ==================================================================
\definecolor{vscodeBlue}{HTML}{569CD6}
\definecolor{vscodeOrange}{HTML}{CE9178}
\definecolor{vscodeGreen}{HTML}{6A9955}
\definecolor{vscodePurple}{HTML}{C586C0}
\definecolor{vscodeGray}{HTML}{9B9B9B}
\definecolor{codeBackground}{HTML}{F8F8F8} % Fond gris très clair
\definecolor{codeText}{HTML}{242424}       % Texte principal (presque noir)
\definecolor{codeGray}{HTML}{A0A0A0}       % Numéros de ligne (gris moyen)

% ==================================================================
% --- MODIFIÉ : CONFIGURATION DU STYLE LISTINGS (LIGHT) ---
% ==================================================================
\lstdefinestyle{vscode}{
    language=Python,
    backgroundcolor=\color{codeBackground},     % Utilise le nouveau fond F8F8F8
    basicstyle=\ttfamily\small\color{codeText}, % Texte principal noir
    keywordstyle=\color{vscodeBlue},
    stringstyle=\color{vscodeOrange},
    commentstyle=\color{vscodeGreen},
    numberstyle=\tiny\color{codeGray},         % Numéros de ligne en gris
    otherkeywords={self, True, False, None},
    keywordstyle=[2]\color{vscodePurple},
    showstringspaces=false,
    breaklines=true,
    frame=none,
    tabsize=4
}

% --- STYLE DE L'OUTPUT ---
\lstdefinestyle{outputstyle}{
    basicstyle=\ttfamily\small\color{codeText}, % Texte principal noir
    breaklines=true,
    frame=none
}

% ==================================================================
% --- MODIFIÉ : DÉFINITION DES CELLULES DE CODE ET OUTPUT (LIGHT) ---
% ==================================================================
% Elles utilisent maintenant le même style "sidebar" que les autres boîtes.

\newtcblisting{codecell}{
  skin=enhanced, % Pour la bordure latérale
  arc=0mm,       % Coins carrés
  boxrule=0pt,   % Pas de cadre
  colback=codeBackground, % Fond clair (F8F8F8)
  borderline west={2pt}{0pt}{vscodeBlue}, % Barre latérale bleue
  fonttitle=\bfseries\color{vscodeBlue}, % Titre en bleu
  listing only,
  listing options={style=vscode, basicstyle=\ttfamily\footnotesize\color{codeText}}, % TEXTE NOIR
  left=3mm, right=3mm, top=2mm, bottom=2mm, % Padding (identique à sidebarstyle)
  boxsep=0mm, % (identique à sidebarstyle)
  breakable  % (identique à sidebarstyle)
}
\newtcblisting{outputcell}{
  skin=enhanced, % Pour la bordure latérale
  arc=0mm,       % Coins carrés
  boxrule=0pt,   % Pas de cadre
  colback=black!5, % Fond gris très clair (Output)
  borderline west={2pt}{0pt}{corrColor}, % Barre latérale grise
  fonttitle=\bfseries\color{corrColor}, % Titre en gris
  listing only,
  listing options={style=outputstyle, basicstyle=\ttfamily\footnotesize\color{codeText}}, % TEXTE NOIR
  left=3mm, right=3mm, top=2mm, bottom=2mm, % Padding (identique à sidebarstyle)
  boxsep=0mm, % (identique à sidebarstyle)
  breakable  % (identique à sidebarstyle)
}

% --- DÉFINITION DES COULEURS POUR DÉF/THÉO/PREUVE/INTUITION/EXEMPLE ---
\definecolor{defColor}{HTML}{1b1f3a}       % Bleu nuit pour les définitions
\definecolor{theoColor}{HTML}{53354a}      % Violet aubergine pour les théorèmes
\definecolor{proofColor}{HTML}{a64942}     % Rouge brique pour les preuves
\definecolor{intuitionColor}{HTML}{16A085} % Turquoise pour l'intuition
\definecolor{exampleColor}{HTML}{4a6982}   % Bleu ardoise pour les exemples
\definecolor{exoColor}{HTML}{1E8449}       % Vert
\definecolor{corrColor}{HTML}{7F8C8D}      % Gris
\definecolor{remarqueColor}{HTML}{D35400} 

% ==================================================================
% --- DÉFINITION DU STYLE DE BASE (MARGES RÉDUITES) ---
% ==================================================================

% --- STYLE "ORGANIQUE" AVEC BARRE LATÉRALE (Pour TOUTES les boîtes) ---
\tcbset{
    sidebarstyle/.style={
        skin=enhanced, % Nécessaire pour les bordures partielles
        arc=0mm,       % Coins carrés
        boxrule=0pt,   % Pas de cadre
        colback=white, % Fond blanc
        colframe=white,
        coltitle=white,
        fontupper=\color{black},
        left=3mm, right=3mm, top=2mm, bottom=2mm, % Padding interne réduit
        boxsep=0mm, % TRÈS IMPORTANT
        breakable
    }
}

% ==================================================================
% --- DÉFINITION DE TOUTES LES CELLULES (TEXTE) ---
% ==================================================================

% --- CELLULES "SIDEBAR" (Pour le contenu théorique) ---
\newtcolorbox{definitionbox}[1][]{
  sidebarstyle,
  borderline west={2pt}{0pt}{defColor}, % Barre latérale gauche
  fonttitle=\bfseries\color{defColor},  % Titre en couleur (sans fond)
  title=Définition\ifstrempty{#1}{}{ : #1}
}
\newtcolorbox{theorembox}[1][]{
  sidebarstyle,
  borderline west={2pt}{0pt}{theoColor},
  fonttitle=\bfseries\color{theoColor},
  title=Théorème\ifstrempty{#1}{}{ : #1}
}
\newtcolorbox{proofbox}[1][]{
  sidebarstyle,
  borderline west={2pt}{0pt}{proofColor},
  fonttitle=\bfseries\color{proofColor},
  title=Preuve\ifstrempty{#1}{}{ : #1}
}
\newtcolorbox{intuitionbox}[1][]{
  sidebarstyle,
  borderline west={2pt}{0pt}{intuitionColor},
  fonttitle=\bfseries\color{intuitionColor},
  title=Intuition\ifstrempty{#1}{}{ : #1}
}
\newtcolorbox{remarquebox}[1][]{
  sidebarstyle,
  borderline west={2pt}{0pt}{remarqueColor},
  fonttitle=\bfseries\color{remarqueColor},
  title=Remarque\ifstrempty{#1}{}{ : #1}
}


% --- CELLULES "SIDEBAR" (Pour les applications) ---
\newtcolorbox{examplebox}[1][]{
  sidebarstyle,
  borderline west={2pt}{0pt}{exampleColor},
  fonttitle=\bfseries\color{exampleColor},
  title=Exemple\ifstrempty{#1}{}{ : #1}
}
\newtcolorbox{exercicebox}[1][]{
    sidebarstyle,
    borderline west={2pt}{0pt}{corrColor},
    fonttitle=\bfseries\color{corrColor},
    title=#1 % <-- On utilise directement l'argument fourni
}
\newtcolorbox{correctionbox}[1][]{
    sidebarstyle,
    borderline west={2pt}{0pt}{exoColor},
    fonttitle=\bfseries\color{exoColor},
    title=#1
}


% =============================================
% --- CORPS DU DOCUMENT ---
% =============================================
\begin{document}

\maketitle

\newpage
\phantomsection
\addcontentsline{toc}{section}{Sommaire}
\tableofcontents 
\newpage

\section{Méthodologies d'évaluation et Arbitrage}

\subsection{La couverture en Delta (Delta Hedging)}

\begin{intuitionbox}[Le principe de base]
Sous certaines hypothèses sur le comportement d'une action, il est possible de déterminer le "juste prix" d'une option. L'idée fondamentale repose sur le fait que la valeur d'une option dépend principalement de deux facteurs :
\begin{itemize}[label=$\cdot$]
    \item Le temps restant avant l'expiration.
    \item Le prix actuel de l'action sous-jacente.
\end{itemize}
Si l'on sait comment le prix de l'option varie en fonction du prix de l'action, on connaît son taux de variation. En termes mathématiques, il s'agit de la dérivée du prix de l'option par rapport au prix de l'action. En finance mathématique, ce taux de variation est appelé le Delta.
\end{intuitionbox}

\begin{definitionbox}[Le Delta ($\Delta$)]
Le Delta d'une option mesure la sensibilité de son prix par rapport aux variations du prix de l'action sous-jacente. 
Par exemple, si le Delta est de $0,5$, cela signifie que pour une augmentation de 1~€ du prix de l'action, le prix de l'option augmentera d'environ 0,50~€.
\end{definitionbox}

\begin{examplebox}[Création d'un portefeuille couvert (sans risque)]
Imaginons un opérateur de marché qui vend une option. Il est exposé au risque si le prix de l'action bouge en sa défaveur. Pour neutraliser ce risque, il va construire le portefeuille suivant :
\begin{itemize}[label=$\cdot$]
    \item Il a vendu une option (position courte : $-1$ option).
    \item Il achète en couverture une quantité d'actions exactement égale au Delta ($\Delta$ actions).
\end{itemize}
Le taux de variation global de ce portefeuille devient alors zéro. Si le prix de l'action varie légèrement, la perte (ou le gain) sur les actions détenues sera exactement compensée par le gain (ou la perte) sur l'option vendue. Mathématiquement, on a éliminé tout le risque du portefeuille.
\end{examplebox}

\begin{remarquebox}[Couverture dynamique et modèle de Black-Scholes]
Le point crucial est que le Delta n'est pas constant : dès que le prix de l'action change, le taux de variation (le Delta) change également. La quantité d'actions détenues dans le portefeuille doit donc être ajustée en permanence. C'est ce qu'on appelle la couverture dynamique (\textit{dynamic hedging}).

Si l'on suppose qu'il est possible de réajuster ce portefeuille \textit{en continu}, on obtient un portefeuille dont la valeur est totalement prévisible et sans risque. Selon les lois de l'arbitrage, un portefeuille sans risque doit rapporter exactement le taux d'intérêt sans risque. C'est cet argument logique strict qui permet de déduire le prix exact de l'option, appelé le prix de Black-Scholes, point de départ de toute la finance mathématique moderne.
\end{remarquebox}

\newpage

\subsection{Exemple pratique et intuitif : Déduction du prix par l'arbitrage}

\begin{examplebox}[Le décor du modèle simplifié binomial]
Pour comprendre la logique fondamentale du \textit{Delta Hedging} sans les mathématiques complexes, utilisons un scénario simple où le temps s'arrête "demain" :
\begin{itemize}[label=$\cdot$]
    \item L'action aujourd'hui : Elle vaut 100~€.
    \item L'action demain (à l'expiration) : Elle ne peut faire que deux choses : monter à 120~€ ou baisser à 80~€.
    \item L'option : Vous vendez une option d'achat (\textit{Call}) donnant le droit d'acheter l'action à 100~€ demain.
    \item Le taux sans risque : La banque (ou l'État) propose un taux d'intérêt garanti de 5\%.
\end{itemize}
\textit{La grande question :} À quel prix exact devez-vous vendre cette option aujourd'hui pour être sûr de ne pas perdre d'argent, ni d'offrir une opportunité d'arbitrage ("d'argent gratuit") au marché ?
\end{examplebox}

\begin{intuitionbox}[Le raisonnement pas à pas]
Étape 1 : Calculer le Delta\\
Que vaudra l'option demain dans nos deux scénarios ?
\begin{itemize}[label=$\cdot$]
    \item Si l'action monte à 120~€ : L'option vaut 20~€ (le droit d'acheter à 100 ce qui vaut 120).
    \item Si l'action baisse à 80~€ : L'option vaut 0~€ (personne n'exerce le droit d'acheter à 100 ce qui vaut 80).
\end{itemize}
L'écart potentiel de l'action est de 40~€ (de 80 à 120). L'écart potentiel de l'option est de 20~€ (de 0 à 20). \\
Le Delta est le rapport entre les deux : $20 / 40 = 0,5$. \\
Le Delta est donc de 0,5. L'option bouge deux fois moins vite que l'action.

\vspace{0.4cm}
Étape 2 : Créer le portefeuille "Sans Risque"\\
Puisque vous avez vendu 1 option (vous êtes en danger si l'action monte), vous vous couvrez en achetant 0,5 action (la valeur du Delta).
Combien d'argent sortez-vous de votre poche aujourd'hui pour construire cela ?
\begin{itemize}[label=$\cdot$]
    \item Achat d'une demi-action à 100~€ : -50~€
    \item Vente de l'option : + Prix de l'option (notre inconnue)
\end{itemize}
Votre investissement net initial aujourd'hui est donc : 50~€ - Prix de l'option.

\vspace{0.4cm}
Étape 3 : Que se passe-t-il demain ? (La "Magie")\\
Regardons la valeur de ce portefeuille dans les deux scénarios finaux :
\begin{itemize}[label=$\cdot$]
    \item Scénario A (L'action explose à 120~€) :
    \begin{itemize}[label=$\triangleright$]
        \item Votre demi-action vaut $0,5 \times 120 = +60$~€.
        \item L'acheteur exerce son option, vous lui devez -20~€.
        \item Valeur totale dans votre poche : $60 - 20 = 40$~€.
    \end{itemize}
    \vspace{0.2cm}
    \item Scénario B (L'action s'effondre à 80~€) :
    \begin{itemize}[label=$\triangleright$]
        \item Votre demi-action vaut $0,5 \times 80 = +40$~€.
        \item L'option expire sans valeur, vous ne devez rien (0~€).
        \item Valeur totale dans votre poche : $40 - 0 = 40$~€.
    \end{itemize}
\end{itemize}
Constat incroyable : Quoi qu'il arrive, vous aurez exactement 40~€ demain. Vous avez créé un investissement garanti à 100\%, c'est-à-dire strictement sans risque.

\vspace{0.4cm}
Étape 4 : Le principe de Non-Arbitrage\\
Puisque vous êtes certain d'avoir 40~€ demain, cet investissement est identique à un placement bancaire garanti à 5\%. \\
Combien devez-vous placer à la banque aujourd'hui pour obtenir 40~€ demain ? \\
La réponse s'obtient par actualisation : $40 / 1,05 \approx 38,10$~€.

La finance moderne repose sur une règle d'or : deux investissements qui rapportent exactement la même chose sans risque doivent coûter le même prix aujourd'hui. La création de votre portefeuille DOIT donc vous coûter exactement 38,10~€.

\vspace{0.4cm}
Conclusion : Déduction du juste prix de l'option\\
Reprenons la formule de notre investissement à l'Étape 2 : \\
Investissement aujourd'hui = 50~€ - Prix de l'option.

Puisque cet investissement doit valoir 38,10~€, il suffit de résoudre l'équation :\\
$50 - \text{Prix de l'option} = 38,10$ \\
Prix de l'option = 11,90~€
\end{intuitionbox}

\begin{remarquebox}[Généralisation aux options en dehors de la monnaie (OTM)]
Cette logique de couverture fonctionne avec absolument n'importe quel Strike (ainsi qu'avec des \textit{Puts}). Puisque le prix de l'option est mathématiquement lié à celui de l'action, il existera toujours un ratio exact (le Delta) permettant d'annuler le risque.

Prenons le même scénario, mais avec un \textit{Call} de Strike 110~€ :
\begin{itemize}[label=$\cdot$]
    \item Nouveau Delta : Si l'action monte à 120~€, l'option vaut 10~€. Si elle baisse à 80~€, l'option vaut 0~€. \\
    Le Delta devient donc $10 / 40 = 0,25$.
    \item Nouveau Portefeuille : Vous vendez l'option et achetez 0,25 action (soit un investissement de 25~€ sur les actions).
    \item La Magie à l'expiration : 
    \begin{itemize}[label=$\triangleright$]
        \item Si 120~€ : Vos 0,25 actions valent 30~€. Vous devez 10~€ à l'acheteur. Reste : 20~€.
        \item Si 80~€ : Vos 0,25 actions valent 20~€. L'option expire à 0~€. Reste : 20~€.
    \end{itemize}
\end{itemize}
Conclusion : Quoi qu'il arrive, la valeur finale est figée à 20~€. Le portefeuille est toujours parfaitement couvert et sans risque. On peut alors appliquer la même logique d'actualisation à 5\% pour déduire le prix exact de ce Call à 110~€.
\end{remarquebox}

\newpage

\subsection{Évaluation des contrats Forward et Coût de Détention}

\begin{theorembox}[Valeur d'un contrat Forward]
Si un actif liquide s'échange aujourd'hui au prix $S_0$, avec un taux de dividende continu $d$ et un taux d'intérêt sans risque composé en continu $r$, alors la valeur $V$ d'un contrat forward pour acheter cet actif au prix $K$ à l'échéance $T$ est :
\begin{equation*}
    V = e^{-rT} (S_0 e^{(r-d)T} - K)
\end{equation*}
Le contrat a une valeur nulle à l'origine si et seulement si le prix d'exercice est fixé au prix forward théorique :
\begin{equation*}
    K = S_0 e^{(r-d)T}
\end{equation*}
\end{theorembox}

\begin{intuitionbox}[Logique du financement par emprunt]
Pour déterminer le juste prix d'un contrat financier, on utilise un raisonnement d'absence d'opportunité d'arbitrage (AOA). L'idée est de construire une stratégie de réplication qui ne coûte strictement rien aujourd'hui (coût initial de $0$). 

Pour acheter l'action aujourd'hui sans sortir d'argent de notre poche, nous sommes obligés de financer cet achat par un emprunt bancaire. Le mécanisme repose sur trois étapes :
\begin{enumerate}
    \item À l'origine ($t=0$) : On emprunte de l'argent à la banque pour acheter l'actif. Le flux de trésorerie net est nul ($+\text{emprunt} - \text{achat} = 0$).
    \item Pendant la durée ($0$ à $T$) : Les dividendes de l'actif sont réinvestis, tandis que la dette à la banque accumule des intérêts au taux sans risque $r$.
    \item À l'échéance ($T$) : On possède $1$ actif, mais on a une dette à rembourser. Pour fixer le prix Forward juste ($K$), il faut que le prix de vente de l'actif couvre exactement le remboursement de cette dette, laissant un gain final de zéro.
\end{enumerate}
\end{intuitionbox}

\begin{proofbox}[L'argument de non-arbitrage]
Pour prouver que le prix forward est $S_0 e^{(r-d)T}$, on construit un portefeuille à coût zéro aujourd'hui :
\begin{itemize}
    \item On emprunte la somme $e^{-dT} S_0$.
    \item On achète une fraction d'actif égale à $e^{-dT}$.
\end{itemize}
À l'échéance $T$ :
\begin{itemize}
    \item La fraction d'actif $e^{-dT}$ a reçu des dividendes qui, une fois réinvestis, nous permettent de détenir exactement $1$ actif complet.
    \item La dette a accumulé des intérêts et vaut désormais $e^{-dT} S_0 \times e^{rT} = S_0 e^{(r-d)T}$.
\end{itemize}
Si le contrat forward nous permet de vendre cet actif au prix $K = S_0 e^{(r-d)T}$, nous recevons exactement la somme nécessaire pour rembourser la dette. Toute différence de prix créerait un profit sans risque (arbitrage).
\end{proofbox}

\begin{examplebox}[Exemple pratique]
Considérons une action $S_0 = 100$, avec $r = 0,05$ et $d = 0,02$ sur une durée $T = 1$.

1. Calcul du prix forward juste :
\begin{equation*}
    F = 100 \times e^{(0,05 - 0,02) \times 1} = 100 \times e^{0,03} \approx 103,05
\end{equation*}

2. Valeur du contrat :
Si vous avez signé un contrat pour acheter cette action à $K = 100$ dans un an :
\begin{itemize}
    \item À l'échéance, vous paierez $100$ pour un actif qui en vaut théoriquement $103,05$.
    \item La valeur de ce contrat aujourd'hui est le gain futur actualisé : 
    $V = e^{-0,05} \times (103,05 - 100) \approx 2,90$.
\end{itemize}
\end{examplebox}

\newpage

\subsection{Parité Call-Put et bornes sur les prix d'options}

\begin{theorembox}[Parité Call-Put]
Si une option d'achat (Call) de prix $C$, une option de vente (Put) de prix $P$ et un contrat forward de valeur $V$ possèdent le même prix d'exercice $K$ et la même échéance $T$, alors :
\begin{equation*}
    C - P = V
\end{equation*}
En utilisant la définition de la valeur d'un contrat forward, nous avons :
\begin{equation*}
    C - P = e^{-rT} (S_0 e^{(r-d)T} - K)
\end{equation*}
\end{theorembox}

\begin{definitionbox}[États d'une option]
Une option est qualifiée selon la position du prix de l'actif par rapport au prix d'exercice $K$ :
\begin{itemize}
    \item Dans la monnaie (In-the-money) : l'option aurait une valeur positive si elle était exercée à l'échéance avec le prix actuel.
    \item En dehors de la monnaie (Out-of-the-money) : l'option n'aurait aucune valeur dans les mêmes conditions.
    \item À la monnaie (At-the-money) : le prix de l'actif est égal au prix d'exercice.
\end{itemize}
Il existe également le concept d'At-the-forward, où le prix d'exercice $K$ est égal au prix forward de l'actif. À ce point précis, la valeur du contrat forward est de $0$, ce qui implique que $C = P$.
\end{definitionbox}

\begin{intuitionbox}[Indépendance vis-à-vis des anticipations]
Un résultat fondamental de la finance mathématique est que nos vues subjectives sur la valeur future moyenne de l'actif n'affectent pas le prix des options. Si nous pensons que le prix a plus de chances de finir au-dessus du prix forward, nous pourrions être tentés de dire que le Call devrait valoir plus que le Put. Cependant, la parité Call-Put montre que si $K$ est égal au prix forward, le Call et le Put doivent avoir exactement la même valeur pour éviter tout arbitrage, quelles que soient les croyances sur la direction du marché.
\end{intuitionbox}

\begin{theorembox}[Bornes sur le prix d'un Call]
Pour une action ne versant pas de dividendes ($d = 0$), le prix d'un Call européen $C_0$ est encadré par le prix actuel de l'action $S_0$ et sa valeur intrinsèque actualisée. Soit $Z_0 = e^{-rT}$ le prix d'un zéro-coupon :
\begin{equation*}
    S_0 > C_0 > S_0 - K Z_0
\end{equation*}
Si les taux d'intérêt sont non négatifs ($r \geq 0$), alors $Z_0 \leq 1$, ce qui implique que le Call vaut toujours plus que sa valeur d'exercice immédiate : $C_0 > S_0 - K$.
\end{theorembox}

\begin{proofbox}[Preuve des bornes]
La borne supérieure $S_0 > C_0$ découle du fait qu'à l'échéance, l'option vaut $\max(S_T - K, 0)$, ce qui est toujours inférieur ou égal au prix de l'action $S_T$. Par le théorème de monotonicité, cette relation doit être vraie à $t = 0$.

Pour la borne inférieure, considérons un portefeuille composé d'une option Call et de $K$ obligations zéro-coupon. À l'échéance $T$ :
\begin{itemize}
    \item Si $S_T > K$ : on exerce l'option pour $K$ et on possède l'action qui vaut $S_T$.
    \item Si $S_T \leq K$ : l'option vaut $0$ et on possède les $K$ de l'obligation.
\end{itemize}
La valeur finale est $\max(S_T, K)$, ce qui est toujours supérieur ou égal à $S_T$. À l'origine, le portefeuille $C_0 + K Z_0$ doit donc valoir plus que l'action $S_0$. En réarrangeant, on obtient $C_0 > S_0 - K Z_0$.
\end{proofbox}
\newpage

\section{Arbres Binomiaux et Évaluation d'Options}

\subsection{Le Modèle à une Période (L'Univers à Deux États)}

\begin{intuitionbox}[Le concept de l'arbre]
Pour évaluer une option classique (vanille), une méthode redoutablement efficace consiste à modéliser les mouvements de l'actif sous-jacent à l'aide d'un arbre. Nous commençons par une situation hautement stylisée (un seul pas de temps, deux états futurs possibles) afin de comprendre la mécanique fondamentale de l'absence d'arbitrage, avant de généraliser.
\end{intuitionbox}

\begin{definitionbox}[Le cadre du modèle binomial simple]
Considérons un univers simplifié avec les paramètres suivants :
\begin{itemize}[label=$\rightarrow$]
    \item \textbf{Taux d'intérêt} : Pour simplifier, nous supposons qu'il est nul ($r = 0$).
    \item \textbf{Actif sous-jacent ($S$)} : Son prix actuel est $S_0 = 100$. Demain, il ne peut prendre que deux valeurs :
    \begin{itemize}[label=$\cdot$]
        \item État "Up" (Hausse) : $S_u = 110$
        \item État "Down" (Baisse) : $S_d = 90$
    \end{itemize}
    \item \textbf{Option d'achat (Call)} : Prix d'exercice $K = 100$.
    \begin{itemize}[label=$\cdot$]
        \item Flux en cas de Hausse : $\max(110 - 100, 0) = 10$
        \item Flux en cas de Baisse : $\max(90 - 100, 0) = 0$
    \end{itemize}
\end{itemize}
\end{definitionbox}

\subsection{Tarification par Couverture (Hedging)}

\begin{intuitionbox}[L'élimination de l'incertitude]
En tant que vendeur de cette option, nous sommes exposés à un risque : perdre 10 si l'état "Up" se réalise, et ne rien perdre dans l'état "Down". La stratégie de couverture consiste à construire un portefeuille aujourd'hui (composé de l'actif sous-jacent et de l'option) dont la valeur de demain sera strictement identique, quel que soit l'état du monde qui se réalise.
\end{intuitionbox}

\begin{proofbox}[Détermination du prix sans arbitrage par la couverture]
Sans connaître l'avenir, nous achetons une quantité fixe $\delta$ d'actions aujourd'hui. Soit $Opt_0$ la valeur de l'option aujourd'hui.
\begin{itemize}[label=$\rightarrow$]
    \item \textbf{Valeur du portefeuille aujourd'hui} : $\Pi_0 = 100\delta - Opt_0$
\end{itemize}
Regardons la valeur de ce portefeuille demain dans les deux états :
\begin{itemize}[label=$\rightarrow$]
    \item Dans l'état "Up" : $\Pi_u = 110\delta - 10$
    \item Dans l'état "Down" : $\Pi_d = 90\delta$
\end{itemize}

\textbf{Étape 1 : Couverture du risque (Delta)}\\
Pour éliminer toute incertitude, le portefeuille doit valoir la même chose dans les deux états ($\Pi_u = \Pi_d$) :
\begin{equation*}
    110\delta - 10 = 90\delta \implies 20\delta = 10 \implies \delta = \frac{1}{2}
\end{equation*}
En détenant une demi-action, notre portefeuille vaudra demain exactement $90 \times (0,5) = 45$, quoi qu'il arrive.

\textbf{Étape 2 : Application du principe de non-arbitrage}\\
Puisque le taux d'intérêt est nul ($r=0$) et que notre portefeuille est sans risque, sa valeur aujourd'hui doit être exactement égale à sa valeur de demain :
\begin{equation*}
    \Pi_0 = 45 \implies 100(0,5) - Opt_0 = 45 \implies 50 - Opt_0 = 45 \implies Opt_0 = 5
\end{equation*}
Le prix unique et sans arbitrage de l'option est de $5$.
\end{proofbox}

\begin{remarquebox}[L'absence des probabilités]
Le point le plus remarquable de ce raisonnement est que les probabilités n'interviennent nulle part. Peu importe qu'il y ait 90\% ou 10\% de chances que l'action monte : l'argument de couverture reste parfaitement valide. La seule condition imposée aux probabilités réelles est qu'elles ne soient ni $0$ ni $1$ (car une probabilité de $1$ de voir l'action monter plus vite que le taux sans risque créerait une opportunité d'arbitrage immédiate).
\end{remarquebox}

\newpage

\subsection{L'Évaluation Risque-Neutre}

\begin{intuitionbox}[L'approche probabiliste]
Que se passe-t-il si nous essayons d'évaluer l'option en utilisant une espérance mathématique ? Pour que cette approche donne le même résultat que notre méthode de couverture (qui est infaillible), nous devons inventer un monde théorique où les investisseurs sont neutres au risque. Dans ce monde, le rendement attendu de l'action est égal au taux sans risque.
\end{intuitionbox}

\begin{theorembox}[Probabilité Risque-Neutre]
Pour retrouver le prix d'arbitrage, on définit une probabilité fictive $p$ (pour l'état "Up") et $1-p$ (pour l'état "Down"), telle que l'espérance de la valeur future de l'action soit exactement égale à sa valeur actuelle (car $r=0$) :
\begin{equation*}
    \mathbb{E}[S_1] = S_0
\end{equation*}
\end{theorembox}

\begin{examplebox}[Calcul du prix via la probabilité risque-neutre]
\textbf{Étape 1 : Trouver la probabilité $p$}\\
Posons l'équation d'espérance de l'action :
\begin{equation*}
    110p + 90(1-p) = 100 \implies 20p + 90 = 100 \implies 20p = 10 \implies p = \frac{1}{2}
\end{equation*}

\textbf{Étape 2 : Évaluer l'option}\\
Dans ce monde risque-neutre, la valeur de l'option aujourd'hui est simplement l'espérance mathématique de ses flux futurs :
\begin{equation*}
    Opt_0 = \mathbb{E}[Flux] = p \times 10 + (1-p) \times 0 = \frac{1}{2}(10) + \frac{1}{2}(0) = 5
\end{equation*}
Nous retrouvons exactement le prix d'arbitrage calculé précédemment, mais avec un calcul beaucoup plus direct.
\end{examplebox}

\subsection{Généralisation et Théorie des Martingales}

\begin{examplebox}[Généralisation algébrique]
Considérons un produit dérivé générique payant $a + \beta$ dans l'état "Up" et $a$ dans l'état "Down".
\begin{itemize}[label=$\rightarrow$]
    \item \textbf{Par couverture} : On égalise les états pour trouver $\delta$ : $110\delta - (a + \beta) = 90\delta - a \implies \delta = \frac{\beta}{20}$. La valeur future garantie est $90(\frac{\beta}{20}) - a = 4,5\beta - a$. L'équation de non-arbitrage donne : $100(\frac{\beta}{20}) - Opt_0 = 4,5\beta - a \implies Opt_0 = a + 0,5\beta$.
    \item \textbf{Par risque-neutre} : L'espérance avec $p=0,5$ donne immédiatement $\mathbb{E}[Opt] = 0,5(a + \beta) + 0,5(a) = a + 0,5\beta$.
\end{itemize}
L'approche risque-neutre est considérablement plus simple à manipuler mathématiquement.
\end{examplebox}

\begin{theorembox}[Propriété de Martingale et Absence d'Arbitrage]
Pourquoi la tarification risque-neutre fonctionne-t-elle ? Une fois la probabilité $p$ fixée de manière à ce que $\mathbb{E}[S] = S_0$, l'opérateur espérance étant linéaire, cette relation s'applique à toute combinaison d'instruments sur le marché.
Si l'on note $B$ l'actif sans risque, on obtient le système suivant :
\begin{align*}
    \mathbb{E}[B_1] &= B_0 \\
    \mathbb{E}[S_1] &= S_0 \\
    \mathbb{E}[Opt_1] &= Opt_0
\end{align*}
La valeur actuelle de \textbf{tout} instrument est l'espérance de sa valeur future. Si un portefeuille vaut zéro aujourd'hui, son espérance future est zéro. Il est donc impossible que ce portefeuille soit toujours positif sans pouvoir être négatif : l'arbitrage est structurellement impossible sous cette probabilité.
\end{theorembox}

\begin{remarquebox}[Unicité du prix]
Les deux arguments (couverture et risque-neutre) fonctionnent en sens inverse :
\begin{itemize}[label=$\rightarrow$]
    \item L'argument risque-neutre montre qu'un prix précis est exempt d'arbitrage (il fixe une borne inférieure).
    \item L'argument par couverture montre que tous les autres prix peuvent être arbitrés (il fixe une borne supérieure).
\end{itemize}
Dans notre modèle binomial, ces deux bornes coïncident parfaitement. Le marché est dit "complet", et nous sommes laissés avec un prix d'arbitrage unique.
\end{remarquebox}

\subsection{La Tarification par Réplication}

\begin{intuitionbox}[La Loi du Prix Unique et la Réplication]
Il existe une troisième manière d'interpréter le prix d'arbitrage d'une option. Au lieu de vendre l'option et de nous couvrir, nous pouvons chercher à fabriquer un clone financier de cette option.
Si nous parvenons à construire un portefeuille composé d'actions et d'emprunts qui génère strictement les mêmes flux que l'option dans tous les états futurs possibles, le principe d'absence d'arbitrage impose que le coût de construction de ce portefeuille aujourd'hui soit exactement égal au prix de l'option.
\end{intuitionbox}

\begin{remarquebox}[L'argument géométrique]
Pourquoi est-il toujours possible de cloner l'option dans notre modèle ? 
La valeur d'un portefeuille d'actions et d'obligations à l'échéance est une fonction affine du prix de l'action, soit Valeur = Delta fois le prix de l'action + Montant obligataire. Graphiquement, c'est une ligne droite.
L'option définit deux points cibles à atteindre : son flux en cas de hausse, et son flux en cas de baisse. Or, par deux points, il passe une et une seule droite. Il existe donc une combinaison unique et exacte d'actions et d'emprunt qui réplique parfaitement l'option.
\end{remarquebox}

\begin{examplebox}[Réplication d'un Call sur Apple Inc.]
Imaginons que nous voulions évaluer une option d'achat (Call) sur l'action Apple (ticker : AAPL).
\begin{itemize}[label=$\rightarrow$]
    \item L'action AAPL cote aujourd'hui 200 \$.
    \item Demain, à la clôture, suite à une annonce de résultats, l'action sera soit à 220 \$ (Hausse), soit à 180 \$ (Baisse).
    \item Le taux d'intérêt à un jour d'un bon du Trésor américain (US Treasury Bill, ticker : BIL) est supposé nul (0\%) pour simplifier.
    \item Nous étudions un Call AAPL de prix d'exercice (Strike) à 200 \$.
\end{itemize}

Étape 1 : Identifier les flux cibles de l'option
\begin{itemize}[label=$\cdot$]
    \item Si AAPL monte à 220 \$, le Call vaut 20 \$.
    \item Si AAPL baisse à 180 \$, le Call vaut 0 \$.
\end{itemize}

Étape 2 : Construire le portefeuille de réplication
Nous cherchons à détenir une quantité Delta d'actions AAPL et un montant B en obligations US Treasury. Nous posons notre système d'équations pour le lendemain :
\begin{itemize}[label=$\cdot$]
    \item Scénario Hausse : 220 Delta + B = 20
    \item Scénario Baisse : 180 Delta + B = 0
\end{itemize}

En soustrayant la deuxième équation à la première, on obtient : 
40 Delta = 20, ce qui donne Delta = 0,5.
En remplaçant Delta dans la deuxième équation : 
180(0,5) + B = 0, ce qui donne 90 + B = 0, donc B = -90.

Étape 3 : Interprétation financière et tarification
Le calcul mathématique nous dicte la stratégie exacte à suivre aujourd'hui :
\begin{itemize}[label=$\cdot$]
    \item Acheter 0,5 action AAPL.
    \item Vendre à découvert pour 90 \$ de bons du Trésor américain (ce qui équivaut à emprunter 90 \$ à la banque).
\end{itemize}

Vérifions ce clone demain : si AAPL vaut 180 \$, notre demi-action vaut 90 \$, ce qui sert tout juste à rembourser notre dette de 90 \$ au Trésor. Il nous reste 0 \$ (exactement comme le Call). Si AAPL vaut 220 \$, notre demi-action vaut 110 \$, nous remboursons les 90 \$, il nous reste 20 \$ net en poche (exactement comme le Call).

Quel est le coût de création de ce clone aujourd'hui ?
Coût net = Achat des actions AAPL - Argent emprunté via les Treasury Bills
Coût net = 0,5 fois 200 \$ - 90 \$ = 100 \$ - 90 \$ = 10 \$.

Le clone parfait coûte 10 \$. Par conséquent, le prix d'arbitrage de l'option Call AAPL sur le marché doit être très exactement de 10 \$.
\end{examplebox}

\begin{theorembox}[Symétrie de la réplication]
Sur un marché complet, l'action, l'obligation sans risque et l'option sont structurellement liés. L'équation de valorisation peut être manipulée algébriquement pour répliquer n'importe quel instrument à partir des deux autres :
\begin{itemize}[label=$\rightarrow$]
    \item Hedging : Action possédée + Option vendue permet de générer des flux identiques à une Obligation sans risque.
    \item Réplication : Action possédée + Obligation empruntée permet de générer des flux identiques à une Option d'achat.
    \item Synthèse d'actif : Option achetée + Obligation détenue permet de générer des flux identiques à une Action.
\end{itemize}
\end{theorembox}

\subsection{Les Limites du Modèle : Le Marché Incomplet}

\begin{intuitionbox}[L'ajout d'un troisième état]
Le modèle binomial à deux états est parfait mathématiquement, mais irréaliste économiquement. Pour s'approcher de la réalité, supposons un instant que l'action (actuellement à 100) puisse prendre trois valeurs demain : monter à 110 (A), stagner à 100 (C), ou baisser à 90 (B). Le Call a toujours un Strike de 100.
Essayons d'appliquer nos méthodes de valorisation à cet univers à trois états.
\end{intuitionbox}

\begin{proofbox}[L'échec de la couverture et l'apparition d'inégalités]
Tentons de nous couvrir en vendant l'option et en achetant une quantité y d'actions. Notre portefeuille demain vaudra :
\begin{itemize}[label=$\cdot$]
    \item État A (110) : $110y - 10$
    \item État C (100) : $100y$
    \item État B (90) : $90y$
\end{itemize}
Pour éliminer le risque, ces trois montants doivent être égaux. Mathématiquement, c'est impossible : nous avons 3 équations simultanées à résoudre mais une seule variable (y).
Si nous utilisons le ratio y = 0,5 du modèle précédent, notre portefeuille vaudra 45 (en A), 50 (en C) ou 45 (en B). Le risque n'est pas éliminé, mais nous savons que le portefeuille vaudra au minimum 45. Par conséquent, sa valeur initiale doit être d'au moins 45 :
\begin{equation*}
    100(0,5) - Opt_0 \ge 45 \implies Opt_0 \le 5
\end{equation*}
L'argument d'arbitrage ne donne plus un prix exact, mais une borne supérieure.
\end{proofbox}

\begin{examplebox}[L'échec de la probabilité risque-neutre unique]
Cherchons une probabilité risque-neutre telle que l'espérance de l'action demain soit 100.
Soit $P_A, P_B, P_C$ les probabilités des trois états.
\begin{equation*}
    110 P_A + 100 P_C + 90 P_B = 100
\end{equation*}
Sachant que $P_C = 1 - P_A - P_B$, l'équation se simplifie pour donner : $P_A = P_B$.
Posons $P_A = P_B = p$. La probabilité $p$ peut être n'importe quel chiffre compris entre 0 et 0,5. Il existe donc une infinité de probabilités risque-neutres valides.
Le prix de l'option étant son espérance, nous avons :
\begin{equation*}
    Opt_0 = 10 \times p + 0 \times P_C + 0 \times P_B = 10p
\end{equation*}
Puisque p est compris entre 0 et 0,5, le prix du Call est un intervalle compris entre 0 et 5.
\end{examplebox}

\begin{theorembox}[Définition d'un Marché Incomplet]
Le modèle à trois états illustre le concept fondamental de marché incomplet. Un marché est incomplet lorsqu'il y a plus d'états futurs du monde que d'actifs financiers indépendants permettant de s'y couvrir.
Dans un tel marché :
\begin{itemize}[label=$\rightarrow$]
    \item La réplication parfaite d'un produit dérivé est impossible.
    \item L'absence d'opportunité d'arbitrage ne suffit plus à fixer un prix unique.
    \item Les mathématiques ne fournissent qu'un intervalle de prix exempts d'arbitrage (ici [0, 5]).
\end{itemize}
Le prix exact qui sera négocié sur le marché ne dépendra plus d'une équation, mais des préférences au risque et de la psychologie des acteurs du marché.
\end{theorembox}

\subsection{Modèle à Plusieurs Périodes}

\begin{intuitionbox}[L'astuce du temps continu pour sauver le modèle]
Nous venons de voir qu'ajouter un troisième état final (stagnation à 100) détruit l'unicité du prix et rend le marché incomplet. 
Pour modéliser des mouvements de prix réalistes sans perdre la puissance mathématique de l'arbitrage, la solution consiste à découper le temps. Si nous supposons que les prix évoluent de manière continue, l'action ne fait pas un grand saut en une journée, elle fait plusieurs petits sauts successifs.
\end{intuitionbox}

\newpage

\begin{examplebox}[Construction d'un arbre binomial à deux pas]
Plutôt que d'envisager une variation de plus ou moins 10 sur une journée entière, nous découpons la journée en deux demi-journées, avec des variations de plus ou moins 5.
\begin{itemize}[label=$\rightarrow$]
    \item Instant initial (t = 0) : L'action vaut 100.
    \item Étape intermédiaire (t = 0,5) : L'action fait face à un univers à deux états. Elle monte à 105 (Nœud D) ou baisse à 95 (Nœud E).
    \item Étape finale (t = 1) : À partir de chaque nœud intermédiaire, l'action fait de nouveau face à deux états.
    \begin{itemize}[label=$\cdot$]
        \item Depuis 105 : elle monte à 110 (État A) ou baisse à 100 (État C).
        \item Depuis 95 : elle monte à 100 (État C) ou baisse à 90 (État B).
    \end{itemize}
\end{itemize}
L'état C (100) est atteint de deux manières différentes : une hausse suivie d'une baisse, ou une baisse suivie d'une hausse. C'est ce que l'on appelle un arbre recombinant.
\end{examplebox}

\vspace{0.5cm}
\begin{center}
\begin{tikzpicture}[>=stealth, sloped]
    % Définition des styles pour les noeuds
    \tikzstyle{node_style} = [circle, draw=vscodeBlue, fill=white, thick, minimum size=1cm, inner sep=2pt, font=\small]
    \tikzstyle{time_style} = [font=\small\color{gray}]

    % Les axes de temps (Placés plus haut à y=4.5 pour éviter le chevauchement)
    \node[time_style] at (0, 4.5) {Aujourd'hui ($t=0$)};
    \node[time_style] at (4, 4.5) {Mi-journée ($t=0,5$)};
    \node[time_style] at (8, 4.5) {Demain ($t=1$)};

    % Placement des noeuds
    % t=0
    \node[node_style] (S0) at (0, 0) {100};

    % t=0.5
    \node[node_style] (D) at (4, 1.5) {105};
    \node[node_style] (E) at (4, -1.5) {95};

    % t=1 (Le noeud A est à y=3)
    \node[node_style] (A) at (8, 3) {110};
    \node[node_style] (C) at (8, 0) {100};
    \node[node_style] (B) at (8, -3) {90};

    % Dessin des flèches
    \draw[->, thick, vscodeBlue] (S0) -- (D) node[midway, above] {Hausse};
    \draw[->, thick, vscodeBlue] (S0) -- (E) node[midway, below] {Baisse};

    \draw[->, thick, vscodeBlue] (D) -- (A);
    \draw[->, thick, vscodeBlue] (D) -- (C);
    
    \draw[->, thick, vscodeBlue] (E) -- (C);
    \draw[->, thick, vscodeBlue] (E) -- (B);
    
    % Noms des états à côté des noeuds finaux
    \node[right=0.2cm of A, font=\small] {État A};
    \node[right=0.2cm of C, font=\small] {État C};
    \node[right=0.2cm of B, font=\small] {État B};
    
    % Noms des états intermédiaires
    \node[above=0.2cm of D, font=\small] {Nœud D};
    \node[below=0.2cm of E, font=\small] {Nœud E};

\end{tikzpicture}
\end{center}
\vspace{0.5cm}

\begin{remarquebox}[Le retour à la complétude du marché]
Ce découpage temporel est une avancée théorique majeure. En regardant l'arbre globalement à t=1, nous avons bien trois états finaux possibles (110, 100, 90). Cependant, si nous nous plaçons à n'importe quel nœud spécifique de l'arbre, il n'y a que deux chemins possibles pour l'étape suivante. 
Le marché est donc localement complet à chaque nœud. Nous pourrons utiliser nos techniques d'évaluation (Couverture ou Risque-Neutre) à chaque intersection pour déduire un prix exact, en travaillant à rebours depuis l'échéance finale jusqu'à aujourd'hui.
\end{remarquebox}

\subsection{Évaluation par Couverture dans un Arbre à Deux Périodes}

\begin{intuitionbox}[Le principe d'induction rétrograde]
Pour calculer la valeur de l'option aujourd'hui, nous devons procéder à l'envers (calcul à rebours). Puisque le marché est complet à chaque nœud, nous pouvons appliquer notre argument de couverture locale. Nous calculons d'abord la valeur de l'option à la mi-journée (nœuds D et E) en utilisant les flux finaux (nœuds A, B et C). Une fois ces valeurs intermédiaires connues, nous recommençons l'opération pour trouver la valeur initiale à t=0.
\end{intuitionbox}

\newpage

\begin{examplebox}[Calcul pas à pas par couverture (Delta Hedging)]
Considérons notre option d'achat (Call) de prix d'exercice K = 100. Son profil de flux à l'échéance (t=1) est : 10 en A, 0 en C, 0 en B. Le taux d'intérêt sans risque est toujours supposé nul.

\begin{enumerate}[label=Étape \arabic* :, leftmargin=*]
    
    \item Évaluation au nœud E (L'action vaut 95 à t=0,5)
    \begin{itemize}[label=$\cdot$]
        \item Si nous sommes au nœud E, l'action finira à l'échéance soit à 100 (État C), soit à 90 (État B). 
        \item Dans les deux cas, le Call de Strike 100 ne vaudra strictement rien (0).
        \item Conclusion immédiate : La valeur de l'option au nœud E est Opt(E) = 0.
    \end{itemize}
    
    \vspace{0.2cm}
    \item Évaluation au nœud D (L'action vaut 105 à t=0,5)
    \begin{itemize}[label=$\cdot$]
        \item À partir de D, l'action peut monter à 110 (le Call vaudra 10) ou baisser à 100 (le Call vaudra 0).
        \item Construisons le portefeuille de couverture : nous achetons une quantité delta d'actions.
        \item Valeur future du portefeuille : 110 delta - 10 (si hausse) et 100 delta - 0 (si baisse).
        \item Pour éliminer le risque, nous égalisons : 110 delta - 10 = 100 delta, ce qui donne 10 delta = 10, soit delta = 1.
        \item Avec 1 action, la valeur garantie à t=1 est de 100.
        \item Par arbitrage, le portefeuille au nœud D vaut 100 : 105(1) - Opt(D) = 100, ce qui implique Opt(D) = 5.
    \end{itemize}

    \vspace{0.2cm}
    \item Évaluation au nœud initial (L'action vaut 100 à t=0)
    \begin{itemize}[label=$\cdot$]
        \item Nous connaissons désormais les valeurs à la mi-journée : 5 si l'action monte à 105, 0 si elle baisse à 95.
        \item Construisons le portefeuille initial : nous achetons une nouvelle quantité delta d'actions.
        \item Valeur future (à t=0,5) : 105 delta - 5 (si hausse vers D) et 95 delta - 0 (si baisse vers E).
        \item Égalisation des flux : 105 delta - 5 = 95 delta, ce qui donne 10 delta = 5, soit delta = 0,5.
        \item Avec 0,5 action, la valeur garantie à t=0,5 est de 95(0,5) = 47,5.
        \item Par arbitrage, la valeur initiale du portefeuille aujourd'hui doit être de 47,5 : 100(0,5) - Opt(0) = 47,5.
        \item Le prix final de l'option aujourd'hui est Opt(0) = 2,5.
    \end{itemize}

\end{enumerate}
\end{examplebox}

\begin{remarquebox}[Couverture dynamique]
Cet exemple démontre mathématiquement la notion de couverture dynamique évoquée précédemment. Remarquez que le delta (la quantité d'actions à détenir pour se couvrir) n'est pas constant. Il vaut 0,5 au début de la journée. Si l'action monte à 105, le vendeur de l'option devra ajuster son portefeuille en achetant 0,5 action supplémentaire pour que son delta passe à 1. La couverture parfaite exige un ajustement à chaque pas de temps.
\end{remarquebox}

\subsection{Évaluation Risque-Neutre dans un Arbre à Deux Périodes}

\begin{intuitionbox}[L'avantage de la méthode probabiliste]
Bien que l'argument de couverture nous garantisse un prix mathématiquement parfait, les calculs de portefeuilles deviennent vite fastidieux. L'approche par la probabilité risque-neutre, bien que plus abstraite, est beaucoup plus simple à appliquer. Elle donnera toujours un résultat strictement identique à la méthode par couverture.
Il existe deux manières de l'utiliser : en remontant l'arbre pas à pas (induction rétrograde), ou en calculant directement les probabilités globales d'atteindre chaque état final.
\end{intuitionbox}

\newpage

\begin{examplebox}[Méthode 1 : Calcul pas à pas (Induction rétrograde)]
Nous reprenons le même arbre. Le taux d'intérêt est nul, donc l'espérance du prix futur de l'action doit égaler son prix actuel à chaque nœud.

\begin{enumerate}[label=Étape \arabic* :, leftmargin=*]
    \item Évaluation au nœud D (L'action vaut 105)
    \begin{itemize}[label=$\cdot$]
        \item L'action peut passer à 110 ou 100. La probabilité p (Hausse) doit vérifier : 110p + 100(1-p) = 105.
        \item La seule probabilité valide est p = 1/2.
        \item La valeur de l'option est son espérance : Opt(D) = 0,5(10) + 0,5(0) = 5.
    \end{itemize}

    \vspace{0.2cm}
    \item Évaluation au nœud E (L'action vaut 95)
    \begin{itemize}[label=$\cdot$]
        \item Les flux finaux de l'option sont nuls dans tous les scénarios (États C et B).
        \item L'espérance est de zéro quelle que soit la probabilité. Opt(E) = 0.
    \end{itemize}

    \vspace{0.2cm}
    \item Évaluation initiale (L'action vaut 100)
    \begin{itemize}[label=$\cdot$]
        \item L'action passe à 105 ou 95. Pour que l'espérance soit 100, la probabilité de Hausse doit de nouveau être p = 1/2.
        \item La valeur initiale de l'option est l'espérance de ses valeurs futures intermédiaires (Opt(D) et Opt(E)).
        \item Opt(0) = 0,5(5) + 0,5(0) = 2,5.
    \end{itemize}
\end{enumerate}
Le résultat est identique au prix d'arbitrage obtenu par la méthode de couverture en delta.
\end{examplebox}

\begin{examplebox}[Méthode 2 : Calcul par les probabilités globales des chemins]
Plutôt que de calculer les valeurs intermédiaires de l'option, nous pouvons calculer directement la probabilité risque-neutre de nous retrouver dans chaque état final à t=1, puis calculer l'espérance finale en une seule fois.
Sachant qu'à chaque nœud la probabilité de Hausse est de 1/2 :

\begin{itemize}[label=$\cdot$]
    \item État A (Action à 110) : Il faut 2 hausses successives. Probabilité = 1/2 fois 1/2 = 1/4.
    \item État C (Action à 100) : Il faut 1 hausse et 1 baisse (ou inversement). Probabilité = 1/4 + 1/4 = 1/2.
    \item État B (Action à 90) : Il faut 2 baisses successives. Probabilité = 1/2 fois 1/2 = 1/4.
\end{itemize}

Calcul direct de l'espérance de l'option aujourd'hui :
Opt(0) = 1/4(10) + 1/2(0) + 1/4(0) = 10/4 = 2,5.

Cette méthode est extrêmement puissante car elle permet de sauter toutes les étapes intermédiaires de calcul.
\end{examplebox}

\begin{theorembox}[Généralisation et Tarification Dynamique]
L'unique caractéristique qualitative cruciale de cet argument est qu'à chaque pas de temps, l'actif ne peut évoluer que vers deux nouvelles valeurs possibles. C'est ce qui garantit la complétude locale du marché.
\begin{itemize}[label=$\rightarrow$]
    \item Le nombre total de pas de temps n'a aucune importance mathématique. On peut subdiviser le temps en autant d'étapes que l'on souhaite (10, 100, 1000 étapes). Tant que l'arbre reste binomial à chaque nœud, le prix d'arbitrage sera unique.
    \item L'interprétation par la réplication reste valide : le prix de l'option (2,5) est exactement la somme d'argent qu'il faut investir aujourd'hui pour fabriquer un clone parfait de l'option.
    \item Ce clone nécessite un rééquilibrage continu. La quantité d'actions à détenir (le Delta) et le montant de cash (emprunt ou placement) changent à chaque nœud de l'arbre en fonction du temps et du prix de l'action. C'est l'essence même de la gestion de portefeuille dynamique.
\end{itemize}
\end{theorembox}

\subsection{Généralisation à N Périodes}

\begin{intuitionbox}[La mécanique de l'arbre global]
Il n'y a rien de magique dans un modèle à deux étapes. Tant que chaque nœud de l'arbre se divise exactement en deux branches, la logique d'induction rétrograde garantit un prix sans arbitrage unique, que l'arbre compte 2, 4 ou 10 000 étapes.
Au lieu de calculer à rebours (nœud par nœud, couche par couche) depuis l'échéance finale jusqu'à aujourd'hui, nous pouvons directement calculer les probabilités globales d'atteindre chaque nœud final de l'arbre, puis faire l'espérance mathématique des flux finaux.
\end{intuitionbox}

\newpage

\begin{examplebox}[Exemple d'un arbre à 4 pas de temps (Taux d'intérêt = 0)]
Imaginons un arbre à 4 étapes, où le cours de l'action monte ou baisse de 2,5 à chaque pas. L'action part de 100 à t=0. 
Puisque le taux d'intérêt est nul, la probabilité risque-neutre de hausse (p) est obligatoirement de 0,5 à chaque nœud pour maintenir l'espérance locale.
Calculons la probabilité risque-neutre globale de chaque état final à la 4ème étape :

\begin{itemize}[label=$\cdot$]
    \item Arrivée à 110 : Il faut exactement 4 hausses successives.
    Il n'y a qu'un seul chemin. Probabilité : $1 \times (1/2)^4 = 1/16$.
    
    \item Arrivée à 105 : Il faut exactement 3 hausses et 1 baisse.
    Le nombre de combinaisons possibles est donné par le coefficient binomial $C_4^1 = 4$.
    Probabilité : $4 \times (1/2)^4 = 4/16 = 1/4$.
    
    \item Arrivée à 100 : Il faut exactement 2 hausses et 2 baisses.
    Nombre de combinaisons $C_4^2 = 6$.
    Probabilité : $6 \times (1/2)^4 = 6/16 = 3/8$.
    
    \item Arrivée à 95 : Il faut exactement 1 hausse et 3 baisses.
    Nombre de combinaisons $C_4^3 = 4$.
    Probabilité : $4 \times (1/2)^4 = 4/16 = 1/4$.
    
    \item Arrivée à 90 : Il faut exactement 0 hausse et 4 baisses.
    Il n'y a qu'un seul chemin. Probabilité : $1 \times (1/2)^4 = 1/16$.
\end{itemize}

Tarification d'un Call de Strike 100 :
Nous multiplions la probabilité de chaque état par la valeur intrinsèque du Call dans cet état.
Opt(K=100) = (1/16) fois (110 - 100) + (1/4) fois (105 - 100) + 0 + 0 + 0
Opt(K=100) = 10/16 + 5/4 = 10/16 + 20/16 = 30/16 = 15/8.
\end{examplebox}

\vspace{0.5cm}
\begin{center}
\begin{tikzpicture}[>=stealth, sloped]
    % Style pour les mini noeuds d'un grand arbre
    \tikzstyle{mini_node} = [circle, draw=vscodeBlue, fill=codeBackground, thick, minimum size=0.4cm, inner sep=0pt, font=\tiny]
    \tikzstyle{time_style} = [font=\small\color{gray}]

    % Dimensions de l'arbre
    \def\N{5} % Nombre d'étapes
    \def\dx{1.8} % Espace horizontal entre les étapes
    \def\dy{0.8} % Espace vertical entre les noeuds

    % Lignes de temps
    \foreach \i in {0,...,\N} {
        \node[time_style] at (\i*\dx, \N*\dy*0.5 + 0.8) {$t=\i$};
    }

    % Génération automatique des noeuds et des liens (Arbre Recombinant)
    \foreach \i in {0,...,\N} {
        \foreach \j in {0,...,\i} {
            % Calcul de la position Y : Centré autour de Y=0
            \pgfmathsetmacro{\ypos}{(\i - 2*\j)*\dy*0.5}
            \node[mini_node] (N-\i-\j) at (\i*\dx, \ypos) {};
            
            % Dessiner les liens vers la génération précédente (si on n'est pas à i=0)
            \ifnum\i>0
                \pgfmathtruncatemacro{\prevI}{\i-1} % CORRECTION ICI
                % Lien depuis le noeud supérieur (baisse)
                \ifnum\j>0
                    \pgfmathtruncatemacro{\prevJup}{\j-1} % CORRECTION ICI
                    \draw[->, vscodeBlue!60] (N-\prevI-\prevJup) -- (N-\i-\j);
                \fi
                % Lien depuis le noeud inférieur (hausse)
                \ifnum\j<\i
                    \pgfmathtruncatemacro{\prevJdown}{\j} % CORRECTION ICI
                    \draw[->, vscodeBlue!60] (N-\prevI-\prevJdown) -- (N-\i-\j);
                \fi
            \fi
        }
    }
    
    % Étiquette pour montrer la multiplication des chemins
    \node[right=0.2cm of N-\N-0, font=\scriptsize] {5 Hausses};
    \node[right=0.2cm of N-\N-5, font=\scriptsize] {5 Baisses};
    \node[right=0.2cm of N-\N-2, font=\scriptsize] {Chemins multiples};
    
\end{tikzpicture}
\end{center}
\vspace{0.5cm}

\begin{theorembox}[La Formule Globale pour N pas de temps]
Considérons une action de prix initial S. L'arbre compte $N$ étapes. À chaque étape, l'action monte ou baisse d'un montant absolu $x$.
À la fin de l'arbre, la valeur de l'action dépend uniquement du nombre total de hausses (noté $j$). Si l'action a fait $j$ hausses, elle a obligatoirement fait $(N - j)$ baisses.
Le prix final de l'action $S_j$ dans l'état avec $j$ hausses s'écrit :
\begin{equation*}
    S_j = S + j(x) - (N - j)(x) = S + (2j - N)x
\end{equation*}

Puisque la probabilité d'une hausse est de $1/2$ (avec un taux nul), la probabilité totale d'atteindre ce prix final $S_j$ suit une distribution binomiale :
\begin{equation*}
    Prob(S_j) = \frac{1}{2^N} \binom{N}{j}
\end{equation*}
Où $\binom{N}{j}$ représente le nombre de chemins différents permettant d'obtenir $j$ hausses sur $N$ lancers.

Par conséquent, la valeur actuelle d'un produit dérivé qui génère un flux $f(S)$ à l'échéance est simplement l'espérance de tous les flux terminaux pondérés par cette probabilité :
\begin{equation*}
    Opt_0 = \frac{1}{2^N} \sum_{j=0}^{N} \binom{N}{j} f\left(S + (2j - N)x\right)
\end{equation*}
\end{theorembox}

\newpage

\begin{intuitionbox}[Vers le temps continu et Black-Scholes]
L'objectif ultime de cette modélisation mathématique n'est pas de calculer un arbre géant à la main, mais de comprendre ce qui se passe mathématiquement lorsque l'on fait tendre le nombre d'étapes $N$ vers l'infini. 
Pour que la limite ait un sens mathématique et économique, il faudra réduire la taille de la variation de prix ($x$) au fur et à mesure que le nombre d'étapes ($N$) augmente, tout en préservant certaines propriétés de la volatilité de l'action. 
C'est précisément ce passage à la limite qui nous fera basculer de l'arbre binomial discret au mouvement brownien continu, fondement du modèle de Black-Scholes.
\end{intuitionbox}
\newpage

\section{Les propriétés des options sur actions}

\subsection{Les facteurs influençant le prix des options}

Il existe six facteurs fondamentaux qui influencent le prix des options sur actions :
\begin{enumerate}
    \item Le cours actuel de l'action ($S_0$)
    \item Le prix d'exercice ($K$)
    \item Le temps restant jusqu'à l'échéance ($T$)
    \item La volatilité du cours de l'action ($\sigma$)
    \item Le taux d'intérêt sans risque ($r$)
    \item Les dividendes prévus durant la vie de l'option
\end{enumerate}

Le tableau suivant résume l'effet de l'\textbf{augmentation} d'une seule variable (toutes choses égales par ailleurs) sur le prix des options de type européen :

\renewcommand{\arraystretch}{1.5}
\begin{tabular}{|l|c|c|}
\hline
\textbf{Variable (si augmentation)} & \textbf{Call Européen} & \textbf{Put Européen} \\
\hline
Cours actuel de l'action ($S_0$) & $+$ & $-$ \\
\hline
Prix d'exercice ($K$) & $-$ & $+$ \\
\hline
Temps jusqu'à l'échéance ($T$) & $?$ & $?$ \\
\hline
Volatilité ($\sigma$) & $+$ & $+$ \\
\hline
Taux d'intérêt sans risque ($r$) & $+$ & $-$ \\
\hline
Dividendes & $-$ & $+$ \\
\hline
\end{tabular}

\begin{remarquebox}[Légende]
\begin{itemize}
    \item $+$ : Signifie que le prix de l'option augmente.
    \item $-$ : Signifie que le prix de l'option diminue.
    \item $?$ : L'effet est incertain (dépend de la relation entre le prix spot et le strike).
\end{itemize}
\end{remarquebox}

\begin{intuitionbox}[Impact du taux sans risque]
L'effet des taux d'intérêt sur le prix de l'option est déterminé par la façon dont ils influencent la valeur actuelle du paiement futur (le prix d'exercice $K$) et le coût d'opportunité différé de la transaction.

\begin{itemize}
    \item \textbf{Pour un Call :} L'acheteur a le droit d'acheter (payer $K$) plus tard. Une hausse des taux réduit la valeur actuelle de cette dette future ($Ke^{-rT}$). Payer "moins cher" en valeur actuelle est un avantage $\to$ \textbf{Le prix du Call augmente}.
    \item \textbf{Pour un Put :} L'acheteur a le droit de vendre (recevoir $K$) plus tard. Une hausse des taux réduit la valeur actuelle de cette rentrée d'argent future. C'est un désavantage $\to$ \textbf{Le prix du Put diminue}.
\end{itemize}
\end{intuitionbox}

\subsection{Synthèse de l'Impact des Taux d'Intérêt ($r$)}

L'effet des taux d'intérêt sur le prix d'une option est déterminé par la façon dont ils influencent la valeur actuelle du paiement futur (le prix d'exercice $K$) et le coût d'opportunité de différer une transaction.

\renewcommand{\arraystretch}{1.5}
\begin{tabular}{|p{0.25\linewidth}|p{0.32\linewidth}|p{0.32\linewidth}|}
\hline
\textbf{Caractéristique} & \textbf{Option d'Achat (Call)} $\nearrow$ & \textbf{Option de Vente (Put)} $\searrow$ \\
\hline
\textbf{Droit Acquis} & Droit d'\textbf{Acheter} à $K$ dans le futur. & Droit de \textbf{Vendre} à $K$ dans le futur. \\
\hline
\textbf{Nature de $K$} & Coût futur (paiement à faire). & Gain futur (paiement à recevoir). \\
\hline
\multicolumn{3}{|c|}{\textbf{Si $r$ augmente}} \\
\hline
\textbf{Impact sur $K$ Actualisé} & La valeur actuelle du coût $K$ diminue (actualisation plus forte). & La valeur actuelle du gain $K$ diminue (actualisation plus forte). \\
\hline
\textbf{Coût d'Opportunité} & Avantage de garder son argent et le placer au taux élevé $r$. & Inconvénient d'attendre pour recevoir $K$ (perte de rendement). \\
\hline
\textbf{Conclusion} & Prix du Call ($C$) \textbf{augmente}. & Prix du Put ($P$) \textbf{diminue}. \\
\hline
\end{tabular}

\begin{intuitionbox}[Le Principe de Base]
La clé est de considérer l'option comme un instrument qui vous permet de différer un flux de trésorerie (le prix d'exercice $K$) jusqu'à la date d'échéance ($T$).

\begin{itemize}
    \item \textbf{Différer un Coût (Call) :} Plus les taux d'intérêt sont élevés, plus il est intéressant de retarder un paiement (le $K$ du call), car l'argent que vous gardez entre-temps vous rapporte plus. Cet avantage augmente le prix du call.
    \item \textbf{Différer un Gain (Put) :} Plus les taux d'intérêt sont élevés, plus il est coûteux de retarder la réception d'un gain (le $K$ du put), car vous perdez le rendement élevé que vous auriez pu obtenir en plaçant cet argent immédiatement. Ce désavantage diminue le prix du put.
\end{itemize}

L'effet est toujours le même : les taux d'intérêt élevés sont un \textbf{avantage} pour celui qui doit payer plus tard (l'acheteur de call) et un \textbf{désavantage} pour celui qui doit être payé plus tard (l'acheteur de put).
\end{intuitionbox}

\subsection{Impact des dividendes}

\begin{intuitionbox}[Dividendes et Valeur d'Option]
Les dividendes représentent des sorties de trésorerie pour l'entreprise, ce qui entraîne mécaniquement une baisse du cours de l'action à la date de détachement du coupon. Cette baisse programmée du sous-jacent affecte la valeur des options. L'annonce ou le versement d'un dividende constitue une mauvaise nouvelle pour le détenteur d'une option d'achat (Call), car la probabilité que le cours dépasse le prix d'exercice diminue. À l'inverse, c'est une bonne nouvelle pour le détenteur d'une option de vente (Put), car la baisse du cours augmente ses chances de profit.
\end{intuitionbox}

\subsection{Bornes pour la valeur des options}

Nous cherchons à définir un intervalle de prix rationnel pour les options, en dehors duquel des opportunités d'arbitrage existeraient.

\subsubsection{Borne supérieure du call européen}

L'établissement de la limite haute pour le prix d'une option d'achat est conceptuellement direct. Fondamentalement, une option d'achat est un droit d'acquérir l'actif sous-jacent. Il serait irrationnel que ce simple droit coûte plus cher que l'actif lui-même.

\begin{intuitionbox}[Logique de la borne supérieure]
Si le prix du call $c$ était supérieur au prix de l'action $S_0$, une opportunité d'arbitrage immédiate et triviale existerait. Un investisseur vendrait le call et achèterait l'action, empochant immédiatement la différence positive $c - S_0$. Quelle que soit l'issue à l'échéance, il détient l'action pour couvrir l'éventuel exercice de l'option, conservant ainsi le profit initial sans aucun risque. Personne n'achèterait une option plus cher que l'action, car l'action confère la propriété immédiate et totale, alors que l'option ne confère qu'une propriété potentielle future moyennant un paiement supplémentaire (le strike).
\end{intuitionbox}

Cette intuition se vérifie mathématiquement en comparant les flux finaux. À l'échéance $T$, le gain de l'option est $\max(S_T - K, 0)$. Puisque le prix d'exercice $K$ est positif, ce gain est nécessairement inférieur ou égal au prix de l'action $S_T$. Si l'actif vaut plus que le dérivé à la fin, il doit valoir plus que le dérivé aujourd'hui.

\begin{theorembox}[Borne Supérieure]
La valeur d'une option d'achat européenne ne peut jamais excéder le cours actuel de l'action sous-jacente.
$$ c \le S_0 $$
\end{theorembox}

\subsubsection{Borne inférieure du call européen}

Pour déterminer la valeur minimale d'un call européen sur une action ne versant pas de dividendes, nous employons une méthodologie similaire en comparant deux portefeuilles d'investissement distincts.

\begin{proofbox}[Démonstration par Portefeuilles]
Considérons les deux stratégies suivantes :
\begin{itemize}
    \item \textbf{Portefeuille A :} Une option d'achat européenne ($c$) et une obligation zéro-coupon garantissant un revenu égal à $K$ à l'échéance (valeur actuelle $Ke^{-rT}$).
    \item \textbf{Portefeuille B :} Une action ($S_0$).
\end{itemize}

Analysons la valeur du portefeuille A à l'échéance $T$. L'obligation rapporte $K$.
\begin{itemize}
    \item Si $S_T > K$ : L'option est exercée. Nous utilisons le cash $K$ pour acheter l'action. Valeur finale : $S_T$.
    \item Si $S_T \le K$ : L'option n'est pas exercée. Nous conservons le cash. Valeur finale : $K$.
\end{itemize}
Ainsi, la valeur du portefeuille A à l'échéance est $\max(S_T, K)$.

Puisque $\max(S_T, K)$ est toujours supérieur ou égal à $S_T$ (la valeur du portefeuille B), il s'ensuit par absence d'arbitrage que la valeur actuelle du portefeuille A doit être supérieure ou égale à celle du portefeuille B :
$$ c + Ke^{-rT} \ge S_0 $$
\end{proofbox}

En réarrangeant les termes de cette inégalité, nous obtenons une première borne : $c \ge S_0 - Ke^{-rT}$. Toutefois, comme pour le put, la valeur de l'option ne peut être négative.

\begin{theorembox}[Borne Inférieure du Call Européen]
La valeur d'un call européen est minorée par la différence entre le cours de l'action et la valeur actuelle du prix d'exercice, ou zéro si cette différence est négative :
$$ c \ge \max(S_0 - Ke^{-rT}, 0) $$
\end{theorembox}

\subsubsection{Borne inférieure du put européen}

Pour déterminer la valeur minimale d'un put européen sur une action ne versant pas de dividendes, nous employons une méthodologie similaire à celle du call, en comparant deux portefeuilles d'investissement distincts.

\begin{proofbox}[Démonstration par Portefeuilles]
Considérons les deux stratégies suivantes :
\begin{itemize}
    \item \textbf{Portefeuille C :} Une option de vente européenne ($p$) et une action ($S_0$).
    \item \textbf{Portefeuille D :} Une obligation zéro-coupon garantissant un revenu égal à $K$ à l'échéance (valeur actuelle $Ke^{-rT}$).
\end{itemize}

Analysons la valeur du portefeuille C à l'échéance $T$. Si le cours de l'action $S_T$ est inférieur à $K$, l'option est exercée (gain de $K-S_T$) et nous possédons l'action ($S_T$), ce qui totalise $K$. Si l'action est supérieure à $K$, l'option vaut 0 mais nous possédons l'action $S_T$.
Ainsi, la valeur du portefeuille C à l'échéance est $\max(S_T, K)$.

Puisque $\max(S_T, K)$ est toujours supérieur ou égal à $K$ (la valeur du portefeuille D), il s'ensuit par absence d'arbitrage que la valeur actuelle du portefeuille C doit être supérieure ou égale à celle du portefeuille D :
$$ p + S_0 \ge Ke^{-rT} $$
\end{proofbox}

En réarrangeant les termes de cette inégalité, nous obtenons une première borne : $p \ge Ke^{-rT} - S_0$. Toutefois, nous savons également qu'une option, représentant un droit et non une obligation, ne peut jamais avoir une valeur négative.

\begin{theorembox}[Borne Inférieure du Put Européen]
La valeur d'un put européen est minorée par la différence entre la valeur actuelle du prix d'exercice et le cours de l'action, ou zéro si cette différence est négative :
$$ p \ge \max(Ke^{-rT} - S_0, 0) $$
\end{theorembox}

\subsection{La parité Call-Put}

Il existe une relation fondamentale et contraignante entre le prix d'un call européen et celui d'un put européen ayant le même prix d'exercice et la même échéance. Cette relation est connue sous le nom de parité Call-Put.

\begin{theorembox}[Relation de Parité Call-Put]
Pour une action ne versant pas de dividendes, la relation suivante doit toujours être vérifiée en l'absence d'opportunité d'arbitrage :
$$ c + Ke^{-rT} = p + S_0 $$
Autrement dit : \textit{Call + Cash = Put + Action}.
\end{theorembox}

\begin{intuitionbox}[Compréhension de la Parité]
Cette égalité découle directement de la comparaison des portefeuilles que nous avons étudiés précédemment.
Imaginons deux investisseurs :
\begin{itemize}
    \item L'investisseur A détient un \textbf{Call} et de l'argent liquide ($Ke^{-rT}$) pour payer le prix d'exercice.
    \item L'investisseur C détient un \textbf{Put} et l'\textbf{Action}.
\end{itemize}
À l'échéance, quelle que soit l'évolution du marché, ces deux investisseurs se retrouveront avec exactement la même richesse finale, soit $\max(S_T, K)$. Puisque les flux futurs sont identiques, les coûts de constitution de ces portefeuilles aujourd'hui doivent être strictement égaux.

Cette équation est puissante car elle permet de déduire le prix d'un call si l'on connaît celui du put (et inversement), sans même avoir besoin d'un modèle d'évaluation complexe.
\end{intuitionbox}

\subsection{Généralisation : La parité Call-Put avec dividendes}

La relation de parité établie précédemment suppose que l'actif sous-jacent ne génère aucun revenu. Or, dans la réalité, la plupart des actions versent des dividendes. Cette distribution de cash modifie l'équilibre car le détenteur de l'action (dans le portefeuille de droite) perçoit les dividendes, tandis que le détenteur du call (dans le portefeuille de gauche) ne les perçoit pas.

Pour rétablir l'égalité, nous devons ajuster le cours de l'action en soustrayant la valeur actuelle des dividendes à percevoir pendant la durée de vie de l'option. Nous notons $I$ la valeur actuelle de ces dividendes.

\begin{theorembox}[Parité Call-Put avec Dividendes]
Si $I$ représente la valeur actuelle des dividendes versés par l'action avant l'échéance $T$, la relation d'absence d'arbitrage devient :
$$ c + Ke^{-rT} = p + S_0 - I $$
\end{theorembox}

\begin{intuitionbox}[Ajustement du sous-jacent]
La logique est identique à celle des contrats forward. Le cours "efficace" pour l'évaluation de l'option n'est pas le cours brut $S_0$, mais le cours net des sorties de trésorerie prévisibles.
On peut réécrire l'équation ainsi :
$$ c + Ke^{-rT} + I = p + S_0 $$
Cela signifie que pour reproduire parfaitement la possession de l'action (qui donne droit au prix final $S_T$ \textbf{et} aux dividendes), le portefeuille composé du Call et du Cash ($Ke^{-rT}$) doit être complété par une somme $I$ équivalente aux dividendes que l'on souhaite "synthétiser".
\end{intuitionbox}
\section{Les stratégies d'échanges impliquant des options}

Les options peuvent être combinées entre elles ou avec l'actif sous-jacent pour créer des stratégies d'investissement adaptées à différents scénarios de marché (hausse, baisse, stabilité ou volatilité).

\subsection{Stratégies impliquant une option et un actif sous-jacent}

La stratégie la plus emblématique de cette catégorie est la \textbf{vente de call couvert} (ou \textit{covered call}).

\begin{definitionbox}[Covered Call]
Cette stratégie consiste à détenir une position longue sur l'action (achat du sous-jacent) tout en prenant simultanément une position courte sur un call (vente d'une option d'achat) sur ce même actif.
\end{definitionbox}

\begin{intuitionbox}[Mécanisme du Covered Call]
Le terme "couvert" vient du fait que l'investisseur possède déjà l'action qu'il pourrait être obligé de livrer si l'option est exercée par l'acheteur.
\begin{itemize}
    \item \textbf{Avantage :} L'investisseur encaisse la prime de l'option, ce qui améliore son rendement si l'action stagne ou monte modérément. Cette prime offre également une légère protection (un "coussin") en cas de baisse du cours de l'action.
    \item \textbf{Inconvénient :} En échange de la prime reçue, l'investisseur renonce à tout profit potentiel lié à une hausse du cours de l'action au-delà du prix d'exercice ($K$). Son gain est plafonné.
\end{itemize}
\end{intuitionbox}

\begin{center}
\begin{tikzpicture}[scale=1.2]
    % Axes
    \draw[->, >=stealth, thick] (0,-2) -- (0,3) node[above] {P\&L};
    \draw[->, >=stealth, thick] (0,0) -- (6,0) node[right] {$S_T$};
    
    % Composantes (Lignes pointillées)
    \draw[dashed, gray] (0,-2) -- (5,3) node[right] {\tiny Long Stock};
    \draw[dashed, gray] (0,1) -- (4,1) -- (5.5,-0.5) node[right] {\tiny Short Call};
    
    % Courbe de la stratégie (Trait noir épais)
    \draw[ultra thick, black] (0,-1) -- (4,2) -- (5.8,2) node[right] {\bfseries Stratégie};
    
    % Annotations
    \draw[dotted] (4,0) node[below] {$K$} -- (4,2);
    \node[left, font=\tiny] at (0,-1) {$-S_0 + Prime$};
\end{tikzpicture}
\end{center}

\subsection{Les Spreads}

Un spread est une stratégie combinant l'achat et la vente d'options de même type (deux calls ou deux puts) sur le même sous-jacent.

\subsubsection{Le Bull Spread (Spread Haussier)}

Comme son nom l'indique, cette stratégie est utilisée lorsque l'investisseur anticipe une hausse du marché.

\textbf{Construction avec des Calls :}
Un bull spread standard est créé en achetant un call à un prix d'exercice bas ($K_1$) et en vendant un call à un prix d'exercice plus élevé ($K_2$).
Puisque le call $K_1$ (plus proche du cours ou dans la monnaie) est plus cher que le call $K_2$ (hors de la monnaie), cette stratégie nécessite un \textbf{investissement initial} (débit).

\textbf{Profil de gain :}
\begin{itemize}
    \item Si le cours finit sous $K_1$ : Les deux options expirent sans valeur. La perte est limitée à l'investissement initial.
    \item Si le cours finit au-dessus de $K_2$ : Le gain sur le call acheté est plafonné par la perte sur le call vendu. Le profit est maximal et constant.
\end{itemize}

% --- GRAPHIQUE BULL SPREAD (Nettoyé + Composants ajoutés) ---
\begin{center}
\begin{tikzpicture}[scale=1.2]
    % Axes
    \draw[->, >=stealth, thick] (0,-2) -- (0,2.5) node[above] {P\&L};
    \draw[->, >=stealth, thick] (0,0) -- (6,0) node[right] {$S_T$};
    
    % Composantes ajoutées (Visuelles/Illustratives)
    % Long Call K1 (Coût plus élevé, donc départ plus bas)
    \draw[dashed, gray] (0,-1.8) -- (2,-1.8) -- (5.5,1.7) node[right, font=\tiny] {Long Call $K_1$};
    % Short Call K2 (Crédit reçu, donc départ positif)
    \draw[dashed, gray] (0,0.8) -- (5,0.8) -- (5.8,0) node[right, font=\tiny] {Short Call $K_2$};

    % Stratégie
    \draw[ultra thick, black] (0,-1) -- (2,-1) -- (5,1.5) -- (5.8,1.5) node[above right] {\bfseries Stratégie};
    
    % Annotations
    \draw[dotted] (2,0) node[below] {$K_1$} -- (2,-1);
    \draw[dotted] (5,0) node[below] {$K_2$} -- (5,1.5);
    \node[left, font=\tiny] at (0,-1) {Débit Initial};
\end{tikzpicture}
\end{center}

\begin{remarquebox}[Bull Spread avec Puts]
Un bull spread peut également être construit en achetant un put avec un prix d'exercice bas ($K_1$) et en vendant un put avec un prix d'exercice plus élevé ($K_2$). Contrairement à la version avec calls, cette combinaison génère un \textbf{flux initial positif} (crédit), car le put vendu ($K_2$) vaut plus cher que le put acheté ($K_1$). Le profil de risque à l'échéance reste similaire (profit si le marché monte).
\end{remarquebox}

\subsubsection{Le Bear Spread (Spread Baissier)}

Le bear spread est la stratégie miroir, utilisée lorsque l'on anticipe une baisse du cours de l'action.

\textbf{Construction avec des Puts :}
On achète un put à un prix d'exercice élevé ($K_2$) et on vend un put à un prix d'exercice plus bas ($K_1$).
Cette stratégie a un coût initial, mais permet de profiter de la baisse jusqu'au niveau $K_1$.

\textbf{Construction avec des Calls :}
On vend un call à un prix d'exercice bas ($K_1$) et on achète un call à un prix d'exercice élevé ($K_2$).
L'investisseur reçoit une prime nette au départ (crédit). Si le marché baisse, il conserve cette prime comme profit maximal. Si le marché monte, ses pertes sont limitées par le call acheté qui agit comme une assurance.

\subsection{Les Butterfly Spreads}

Le butterfly spread est une stratégie de volatilité neutre. Elle implique des positions sur trois prix d'exercice différents, généralement espacés de manière symétrique : $K_1 < K_2 < K_3$.

\begin{definitionbox}[Construction du Butterfly]
Pour créer un butterfly spread en utilisant des calls :
\begin{itemize}
    \item Achat d'un call au prix d'exercice bas $K_1$.
    \item Achat d'un call au prix d'exercice élevé $K_3$.
    \item Vente de \textbf{deux} calls au prix d'exercice intermédiaire $K_2$ (généralement proche du cours actuel $S_0$).
\end{itemize}
\end{definitionbox}

\begin{intuitionbox}[Profil de gain]
Cette stratégie produit un bénéfice maximal si le cours de l'action à l'échéance reste proche de $K_2$. C'est une stratégie de "stabilité".
\begin{itemize}
    \item Si l'action varie significativement (très bas en dessous de $K_1$ ou très haut au-dessus de $K_3$), la stratégie conduit à une perte faible et limitée (égale au coût de mise en place de la stratégie).
    \item Le graphique de gain ressemble à un triangle (ou aux ailes d'un papillon) avec un pic en $K_2$.
\end{itemize}
\end{intuitionbox}

\begin{center}
\begin{tikzpicture}[scale=1.2]
    \draw[->, >=stealth, thick] (0,-1.5) -- (0,2.5) node[above] {P\&L};
    \draw[->, >=stealth, thick] (0,0) -- (7,0) node[right] {$S_T$};
    
    \draw[blue] (0,-1) -- (1.5,-1) -- (4,1.5); 
    \draw[red] (0,-0.5) -- (5.5,-0.5) -- (7,1); 
    \draw[green] (0,0.7) -- (3.5,0.7) -- (5.5,-2.6);
    
    \draw[thick, black] (0,-0.8) -- (1.5,-0.8) -- (3.5,1.8) -- (5.5,-0.8) -- (7,-0.8) node[right] {\bfseries Stratégie};
    
    \node[above left, font=\small] at (1.5,0) {$K_1$};
    \node[below, font=\small] at (3.5,0) {$K_2$};
    \node[above right, font=\small] at (5.5,0) {$K_3$};

\end{tikzpicture}
\end{center}

\subsection{Le Straddle (Stellage)}

Le Straddle est une stratégie purement orientée vers la volatilité. Elle est utilisée lorsqu'un investisseur s'attend à un mouvement violent du cours de l'action, mais ignore dans quelle direction ce mouvement se produira (par exemple, avant l'annonce de résultats d'entreprise incertains).

\begin{definitionbox}[Construction du Straddle]
Un \textbf{Long Straddle} consiste à acheter simultanément :
\begin{itemize}
    \item Un Call au prix d'exercice $K$.
    \item Un Put au même prix d'exercice $K$.
\end{itemize}
Les deux options doivent avoir la même date d'échéance.
\end{definitionbox}

\begin{intuitionbox}[Profil de Risque/Rendement]
Cette stratégie nécessite un investissement initial important (somme des primes du Call et du Put).
\begin{itemize}
    \item \textbf{Perte maximale :} Limitée au coût des primes. Elle survient si le cours de l'action finit exactement à $K$ (les deux options expirent sans valeur).
    \item \textbf{Gain :} Il est potentiellement illimité. Pour être rentable, l'action doit s'éloigner suffisamment de $K$ (à la hausse ou à la baisse) pour couvrir le coût des deux primes. Le profil de gain forme un "V".
\end{itemize}
\end{intuitionbox}

\begin{center}
\begin{tikzpicture}[scale=1.2]
    % Axes
    \draw[->, >=stealth, thick] (0,-2) -- (0,2.5) node[above] {P\&L};
    \draw[->, >=stealth, thick] (0,0) -- (6,0) node[right] {$S_T$};
    
    % Composantes (Lignes pointillées grises)
    % Call
    \draw[dashed, gray] (0,-0.8) -- (3,-0.8) -- (5.5,1.7) node[right, font=\tiny] {Long Call};
    % Put
    \draw[dashed, gray] (0,2.2) -- (3,-0.8) -- (5.5,-0.8) node[right, font=\tiny] {Long Put};
    
    % Stratégie (Trait noir épais)
    % Somme des deux primes = -1.6 (arbitraire pour le dessin)
    \draw[ultra thick, black] (0,1.4) -- (3,-1.6) -- (5.5,0.9) node[above right] {\bfseries Straddle};
    
    % Annotations
    \draw[dotted] (3,0) node[below] {$K$} -- (3,-1.6);
    \node[left, font=\tiny] at (0,-1.6) {Coût Total};
    
    % Break-even points (approximatifs visuellement)
    \draw[fill=black] (1.4,0) circle (1pt) node[below left, font=\tiny] {$BE_{bas}$};
    \draw[fill=black] (4.6,0) circle (1pt) node[below right, font=\tiny] {$BE_{haut}$};
\end{tikzpicture}
\end{center}

\subsection{Le Strangle}

Le Strangle est une variante du Straddle, également utilisée pour parier sur la volatilité, mais avec un coût d'entrée plus faible.

\begin{definitionbox}[Construction du Strangle]
Un \textbf{Long Strangle} consiste à acheter deux options \textit{hors de la monnaie} (Out-of-the-Money) :
\begin{itemize}
    \item Un Put avec un prix d'exercice bas $K_1$.
    \item Un Call avec un prix d'exercice plus élevé $K_2$ (avec $K_1 < K_2$).
\end{itemize}
\end{definitionbox}

\begin{intuitionbox}[Différence avec le Straddle]
Comme les deux options sont achetées "hors de la monnaie", leurs primes sont moins chères que celles d'un Straddle (qui utilise des options "à la monnaie").
\begin{itemize}
    \item \textbf{Avantage :} Le coût initial est plus faible, donc la perte maximale est réduite.
    \item \textbf{Inconvénient :} Pour être rentable, l'action doit faire un mouvement \textbf{encore plus important} que pour un Straddle. Le cours doit dépasser $K_2$ ou passer sous $K_1$ de manière significative.
\end{itemize}
\end{intuitionbox}

\begin{center}
\begin{tikzpicture}[scale=1.2]
    % Axes
    \draw[->, >=stealth, thick] (0,-1.5) -- (0,2.5) node[above] {P\&L};
    \draw[->, >=stealth, thick] (0,0) -- (6,0) node[right] {$S_T$};
    
    % Composantes (Lignes pointillées grises)
    % Long Put K1
    \draw[dashed, gray] (0,1.5) -- (2,-0.5) -- (5.5,-0.5);
    % Long Call K2
    \draw[dashed, gray] (0,-0.5) -- (4,-0.5) -- (5.5,1);
    
    % Stratégie (Trait noir épais)
    % Coût total = -1.0 (somme des deux primes plus faibles)
    \draw[ultra thick, black] (0,1) -- (2,-1) -- (4,-1) -- (5.5,0.5) node[right] {\bfseries Strangle};
    
    % Annotations
    \draw[dotted] (2,0) node[above] {$K_1$} -- (2,-1);
    \draw[dotted] (4,0) node[above] {$K_2$} -- (4,-1);
    \node[left, font=\tiny] at (0,-1) {Coût Total};
    
    % Zone de perte
    \node[font=\tiny, text=red] at (3,-0.7) {Zone de perte};

\end{tikzpicture}
\end{center}
\newpage

\section{Les arbres binomiaux}

L'approche par arbres binomiaux, souvent associée au modèle de Cox-Ross-Rubinstein, propose une méthode discrète et intuitive pour modéliser la dynamique du cours d'une action et évaluer les options.

Contrairement aux modèles continus qui peuvent sembler abstraits, ce modèle repose sur une idée simple : le temps est découpé en périodes, et à chaque étape, le prix ne peut aller que dans deux directions. C'est cette simplicité qui permet de comprendre la mécanique profonde de l'évaluation d'options. Notons d'ailleurs que lorsque la durée des périodes tend vers 0, ce modèle converge mathématiquement vers la loi log-normale utilisée dans le modèle de Black \& Scholes.

\subsection{Modélisation à une période}

Commençons par le cas le plus élémentaire : une seule période de temps $T$.
Soit une action de prix actuel $S_0$ et une option sur cette action (par exemple un Call) de valeur inconnue $f$.

L'hypothèse fondamentale est que le cours suit une marche aléatoire multiplicative. À l'issue de la période $T$, deux états de la nature sont possibles :
\begin{itemize}
    \item \textbf{État Haut (Up) :} Le prix monte à $S_u = S_0 \cdot u$ (avec $u > 1$).
    \item \textbf{État Bas (Down) :} Le prix descend à $S_d = S_0 \cdot d$ (avec $d < 1$).
\end{itemize}

Les payoffs (flux) de l'option à l'échéance dépendront de ces états : $f_u$ si le marché monte, et $f_d$ s'il baisse.

\begin{center}
\begin{tikzpicture}[grow=right, sloped, level distance=3.5cm, sibling distance=2.5cm]
    \node[left] {\textbf{Aujourd'hui} ($t=0$) : $S_0, f$}
    child {
        node[right] {\textbf{Baisse} : $S_d, f_d$}
        edge from parent
        node[below] {facteur $d$}
    }
    child {
        node[right] {\textbf{Hausse} : $S_u, f_u$}
        edge from parent
        node[above] {facteur $u$}
    };
\end{tikzpicture}
\end{center}

\subsection{Le mécanisme de couverture (Delta-Hedging)}

Comment déterminer la valeur $f$ de l'option aujourd'hui ? L'intuition géniale de ce modèle ne repose pas sur la spéculation de la direction du marché, mais sur la **réplication**.

Nous allons construire un portefeuille fictif $\Pi$ composé de :
\begin{itemize}
    \item Une position longue sur $\Delta$ actions.
    \item Une position courte sur 1 option (vente de l'option).
\end{itemize}

La valeur de ce portefeuille à l'instant $T$ est aléatoire :
$$ \Pi_T = \begin{cases} S_0 u \Delta - f_u & \text{si le marché monte} \\ S_0 d \Delta - f_d & \text{si le marché baisse} \end{cases} $$

L'objectif est de choisir $\Delta$ de manière à **éliminer le risque**. Nous voulons que la valeur du portefeuille soit identique, que le marché monte ou qu'il descende.

\begin{proofbox}[Calcul du Ratio de Couverture $\Delta$]
Pour rendre le portefeuille sans risque, nous égalisons les deux issues possibles :
$$ S_0 u \Delta - f_u = S_0 d \Delta - f_d $$

En regroupant les termes en $\Delta$, nous obtenons :
$$ \Delta (S_0 u - S_0 d) = f_u - f_d $$

Ce qui nous donne la formule du Delta :
$$ \Delta = \frac{f_u - f_d}{S_u - S_d} $$
\end{proofbox}

\textbf{Interprétation économique :} Le ratio $\Delta$ représente la sensibilité du prix de l'option aux variations du sous-jacent. En détenant $\Delta$ actions pour chaque option vendue, les gains sur l'action compensent exactement les pertes sur l'option (et inversement).

\subsection{Dérivation de la formule de valorisation}

Puisque nous avons construit un portefeuille sans risque, il ne doit pas exister d'opportunité d'arbitrage. Par conséquent, ce portefeuille doit rapporter exactement le taux sans risque $r$. Si ce n'était pas le cas, des investisseurs pourraient réaliser un profit infini sans risque.

La valeur actuelle du portefeuille ($\Pi_0$) doit donc être égale à sa valeur future certaine actualisée au taux $r$ :
$$ \underbrace{S_0 \Delta - f}_{\text{Coût initial}} = \underbrace{(S_0 u \Delta - f_u)}_{\text{Valeur future certaine}} \cdot e^{-rT} $$

C'est ici que nous intégrons le calcul complet pour isoler le prix de l'option $f$.

\begin{proofbox}[Détail du calcul algébrique]
Notre but est d'exprimer $f$ en fonction des paramètres connus. Partons de l'équation d'absence d'arbitrage ci-dessus :

\textbf{Étape 1 : Isoler $f$}
$$ f = S_0 \Delta - (S_0 u \Delta - f_u) e^{-rT} $$
Factorisons les termes contenant $\Delta$ :
$$ f = S_0 \Delta (1 - u e^{-rT}) + f_u e^{-rT} $$

\textbf{Étape 2 : Substitution du Delta}
Remplaçons $\Delta$ par sa valeur calculée précédemment $\frac{f_u - f_d}{S_0(u - d)}$ :
$$ f = S_0 \left[ \frac{f_u - f_d}{S_0(u - d)} \right] (1 - u e^{-rT}) + f_u e^{-rT} $$
Les termes $S_0$ s'annulent. Pour simplifier, factorisons l'ensemble par le facteur d'actualisation $e^{-rT}$. Notez que le terme $(1 - u e^{-rT})$ devient $(e^{rT} - u)$ une fois $e^{-rT}$ sorti en facteur :
$$ f = e^{-rT} \left[ \frac{f_u - f_d}{u - d} (e^{rT} - u) + f_u \right] $$

\textbf{Étape 3 : Réarrangement en probabilités}
Développons pour regrouper les termes multipliant $f_u$ et ceux multipliant $f_d$ :
$$ f = e^{-rT} \left[ f_u \left( \frac{e^{rT} - u}{u - d} + 1 \right) + f_d \left( \frac{-(e^{rT} - u)}{u - d} \right) \right] $$

Simplifions le coefficient devant $f_u$ (mise au même dénominateur) :
$$ \frac{e^{rT} - u + u - d}{u - d} = \frac{e^{rT} - d}{u - d} $$

Nous définissons cette quantité comme étant $p$. On constate alors que le coefficient devant $f_d$ est exactement $1-p$.
\end{proofbox}

\subsection{La probabilité risque-neutre}

Le calcul précédent nous permet d'aboutir à une formule d'une élégance remarquable, qui ressemble à une espérance mathématique.

\begin{theorembox}[Formule de valorisation binomiale]
La valeur de l'option est l'espérance actualisée de ses flux futurs sous la probabilité risque-neutre $p$ :
$$ f = e^{-rT} [p f_u + (1 - p) f_d] $$
Où la probabilité $p$ est définie par :
$$ p = \frac{e^{rT} - d}{u - d} $$
\end{theorembox}


\input{sections/test.tex}


\end{document}